\subsection{仿射空间的射影完备化}

设 V 是域 K 上的一个（左）向量空间，考虑 K 上的向量空间 \(K_{s} \times V\) ；射影空间 \(\mathbf{P}(\mathbf{K}_{s} \times \mathbf{V})\) 称为与向量空间 V 典范相伴的射影空间。若 V 的维数为 n，则 \(\mathbf{P}(\mathbf{K}_{s} \times \mathbf{V})\) 有相同的维数 n。考虑 \(K_{s} \times V\) 中的仿射超平面 \(V_{1} = \{1\} \times V\) ，它的方向（第 3 号）是子空间 \(V_{0} = \{0\} \times V\) ；若 \(K_{s} \times V\) 的一条（过 0 的）直线不含于 \(V_{0}\) ，则它包含一个点 \((\alpha, x)\) ，其中 \(\alpha \neq 0\) 且 \(x \in V\) ，因而它也包含点 \(\alpha^{-1}(\alpha, x) = (1, \alpha^{-1}x)\) （属于 \(V_{1}\) ）；其逆是直接的，于是可以看出，在 \(V_{1}\) 的点与 \(\mathbf{K}_s \times \mathbf{V}\) 中不含于 \(\mathbf{V}_0\) 的（过 0 的）直线之间存在一一对应，后者的每一条与 \(\mathbf{V}_1\) 交于且仅交于一点。由此可知，映射 \(x \mapsto \phi(x) = \pi(1, x)\) 是 \(\mathbf{V}\) 到射影空间 \(\mathbf{P}(\mathbf{K}_s \times \mathbf{V})\) 中的一个单射（称为典范单射）； \(\mathbf{V}\) 常与此单射下的像等同。 \(\phi(\mathbf{V})\) 在 \(\mathbf{P}(\mathbf{K}_s \times \mathbf{V})\) 中的补集是射影超平面 \(\mathbf{P}(\mathbf{V}_0)\) ，称为 \(\mathbf{P}(\mathbf{K}_s \times \mathbf{V})\) （或按滥用的说法， \(\mathbf{V}\) 的）的无穷远超平面；它的点也称为 \(\mathbf{P}(\mathbf{K}_s \times \mathbf{V})\) （或 \(\mathbf{V}\) ）的无穷远点。若 \((a_t)\) 是 \(\mathbf{V}\) 的一个基，且在 \(\mathbf{K}_s \times \mathbf{V}\) 中取由元素 \(e_t = (0, a_t)\) 与元素 \(e_\omega = (1, 0)\) 组成的基，则 \(\mathbf{P}(\mathbf{K}_s \times \mathbf{V})\) 中的无穷远点是那些指标 \(\omega\) 的齐次坐标为 0 的点。

设 M 是 V 中的一个仿射线性簇（第 3 号），D 是它的方向；M 的典范像 \(\phi(M)\) （在 \(\mathbf{P}(\mathbf{K}_s \times \mathbf{V})\) 中）含于典范像 \(\overline{\mathbf{M}} = \pi(\mathbf{M}_2)\) ，它是向量子空间 \(\mathbf{M}_2\) （ \(\mathbf{K}_s \times \mathbf{V}\) 中由仿射簇 \(\mathbf{M}_1 = \{1\} \times \mathbf{M}\) （ \(\mathbf{K}_s \times \mathbf{V}\) 中的）所生成的）的典范像。更精确地说，若 \((a_t)\) 是 M 的一个生成 M 的仿射自由系，则诸元素 \((1, a_t)\) 构成 \(\mathbf{M}_2\) 的一个基，从而 \(\overline{\mathbf{M}}\) 正是由 \(\phi(M)\) 生成的射影线性簇；若 M 是有限维的，则 \(\overline{\mathbf{M}}\) 与 M 有相同的维数。 \(\phi(M)\) 在 \(\overline{\mathbf{M}}\) 中的补集是 \(\overline{\mathbf{M}}\) 与无穷远超平面的交，且等于典范像 \(\pi(\mathbf{M}_0)\) ，其中 \(\mathbf{M}_0 = \{0\} \times \mathbf{D}\) 。

反之，设 N 是一个不含于无穷远超平面的射影线性簇，并设 \(\mathbf{R} = \overline{\pi}^{-1}(\mathbf{N})\) ； \(R \cap V_{1}\) 是 \(K \times V\) 的一个仿射线性簇，形如 \(\{1\} \times M\) ，其中 M 是 V 的一个仿射线性簇，并且立即可看出，N 是仿射线性簇 \(\overline{M}\) ，它由 \(\phi(M)\) 生成。

因此，在 V 的仿射线性簇与 \(\mathbf{P}(\mathbf{K}_{s} \times \mathbf{V})\) 中不含于无穷远超平面的射影线性簇之间存在一一对应；V 的两个仿射线性簇平行，必要且充分条件是它们生成的射影线性簇与无穷远超平面有相同的交（这有时表述为：所考虑的两个仿射线性簇具有相同的无穷远点）。

